\begin{explanationbox}
\textbf{一、先看图形：这是一个“从椭圆上的点看两个焦点”的问题。}

在纸上先画椭圆及长轴上的 $F_1,F_2$，再取椭圆上一点 $P$，连接 $PF_1,PF_2$。
站在 $P$ 看向两个焦点，两条视线张开的角就是 $\theta=\angle F_1PF_2$。
我们问：\emph{当 $P$ 沿椭圆边界移动时，这两条视线最多能张开多大？}

请先记住两件始终不变的事：
\[
\underbrace{PF_1+PF_2}_{\text{两条视线的总长度}}=2a,\qquad
\underbrace{F_1F_2}_{\text{两个焦点之间的距离}}=2c.
\]
也就是说，三角形 $F_1PF_2$ 的\textbf{底边固定为 $2c$，另外两边的和固定为 $2a$}。
当 $P$ 改变时，真正发生变化的是两条边的分配：一条长一点，另一条就短一点。

\medskip
\textbf{二、为什么短轴端点的张角最大？先用“等腰三角形”理解。}

把 $P$ 想象成两条视线的起点。若 $P$ 靠近某一个焦点，$PF_1$ 与 $PF_2$ 一长一短，
两条视线的方向就不够“对称”。当 $P$ 位于短轴端点时，
\[
PF_1=PF_2=a,
\]
焦点三角形变成等腰三角形，两条视线关于 $y$ 轴对称，张角达到最大。

这并非只凭观察得到的猜想。可以用\textbf{投影}准确说明：
在点 $P$ 处作角 $\theta$ 的角平分线，再取\emph{垂直于角平分线}的方向。
两条视线在这个方向上的投影长度分别为
$PF_1\sin(\theta/2)$ 和 $PF_2\sin(\theta/2)$，方向相反。
所以，连接两个焦点的线段 $F_1F_2$ 在这个方向上的投影长度恰好是
\[
(PF_1+PF_2)\sin\frac{\theta}{2}
=2a\sin\frac{\theta}{2}.
\]
\textbf{线段的投影长度不可能大于线段本身的长度 $2c$}，因此
\[
2a\sin\frac{\theta}{2}\le2c
\quad\Longrightarrow\quad
\boxed{\sin\frac{\theta}{2}\le\frac ca=e.}
\]
只有当 $F_1F_2$ 恰好垂直于角平分线时，投影长度才能等于线段全长；
此时 $PF_1=PF_2$，即 $P$ 在短轴端点。

\medskip
\textbf{三、回到最容易看懂的特殊图形：短轴端点。}

设 $P=(0,b)$。因为焦点关于 $y$ 轴对称，$PF_1=PF_2=a$。
将等腰三角形沿 $PO$（$O$ 为原点）分成两个直角三角形，则
\[
\sin\frac{\theta_{\max}}2
=\frac{OF_2}{PF_2}
=\frac ca=e.
\]
因此
\[
\boxed{\theta_{\max}=2\arcsin e.}
\]
\textbf{记忆画面：短轴端点、等腰焦点三角形、半角直角三角形；对边 $c$，斜边 $a$。}

\medskip
\textbf{四、用“两个焦半径是否相等”理解最大值的本质。}

记 $r_1=PF_1,r_2=PF_2$。由余弦定理可以得到
\[
\sin^2\frac\theta2=1-\frac{b^2}{r_1r_2}.
\]
这里 $b$ 不变，而 $r_1+r_2=2a$ 不变。
\textbf{两个正数的和固定时，乘积在两数相等时最大}：
\[
r_1r_2\le a^2.
\]
乘积越大，$1-b^2/(r_1r_2)$ 越大，张角也越大。
这就把“等腰时最大”从图像直觉变成了可以计算的事实。

\medskip
\textbf{五、离心率为什么能控制角度？}

离心率 $e=c/a$ 描述焦点间距 $2c$ 占焦半径总长 $2a$ 的比例。
如果保持长半轴 $a$ 不变，让 $c$ 增大，则 $b=\sqrt{a^2-c^2}$ 减小，椭圆变扁；
同时
\[
\theta_{\max}=2\arcsin\frac ca
\]
变大。\textbf{椭圆越扁，两个焦点的最大观测张角越大。}
注意：这一判断说的是\emph{最大张角}，并非任意两个不同椭圆上的点都能直接比较。

\medskip
\textbf{六、数值画面：以 $a=5,b=4,c=3$ 为例。}

此时 $e=3/5$。在短轴端点 $P=(0,4)$，
\[
PF_1=PF_2=5,\qquad
\sin\frac{\theta_{\max}}2=\frac35,
\quad\theta_{\max}\approx73.74^\circ.
\]
当 $P$ 移到长轴端点 $(5,0)$ 时，两条视线沿同一方向，$\theta=0$。
所以从长轴端点走向短轴端点，可以直观地看到张角由 $0$ 增大到最大值。

\medskip
\textbf{七、高考识别信号。}

题目若出现“椭圆上某点看两焦点的夹角”“存在直角或钝角的焦点三角形”
“由焦点张角反求离心率”，优先想到
\[
\boxed{\sin\frac\theta2\le e}
\qquad\text{或}\qquad
\boxed{\theta_{\max}=2\arcsin e}.
\]
\end{explanationbox}
